\documentclass[10pt,a4paper]{article}
\usepackage{fontspec,amsmath,amssymb,graphicx}
\usepackage[margin=18mm]{geometry}
\setmainfont{Times New Roman}
\setlength{\parindent}{0pt}
\setlength{\parskip}{5pt}
\allowdisplaybreaks
\sloppy
\emergencystretch=2em
\begin{document}
\section*{VD\_2\_2\_02 — Lời giải bổ sung}
\subsection*{Kiến thức và trình tự giải}
\textbf{Mục tiêu.} Phân loại và đưa phương trình tuyến tính cấp hai về dạng chính tắc, có giữ đầy đủ các hạng cấp thấp.

\textbf{Bước 1 — Chuẩn hóa hệ số.} Viết


\[
A u_{xx}+2B u_{xy}+C u_{yy}+D u_x+E u_y+K u=f.
\]


Hệ số $B$ bằng một nửa hệ số in trước $u_{xy}$. Tính $\delta=B^2-AC$: dương là hyperbolic, âm là elliptic, bằng $0$ là parabolic nếu phần cấp hai còn khác $0$. Phải xét theo từng miền khi hệ số thay đổi hoặc suy biến.

\textbf{Bước 2 — Chọn biến mới.} Với hyperbolic, giải $A(dy)^2-2Bdxdy+C(dx)^2=0$ để tìm hai bất biến $\xi,\eta$ của hai họ đặc trưng. Với parabolic, chỉ có một họ: lấy bất biến $s$ và một biến $z$ độc lập. Với elliptic, không có hai họ đặc trưng thực; hoàn thành bình phương của phần chính để loại hạng hỗn hợp rồi cân bằng hai hệ số dương. Mọi trường hợp đều kiểm tra Jacobian khác $0$.

\textbf{Bước 3 — Viết công thức thay đạo hàm.} Nếu $u(x,y)=v(\xi,\eta)$ thì


\[
\begin{aligned}
u_x&=\xi_xv_\xi+\eta_xv_\eta,\\
u_{xx}&=\xi_x^2v_{\xi\xi}+2\xi_x\eta_xv_{\xi\eta}
+\eta_x^2v_{\eta\eta}+\xi_{xx}v_\xi+\eta_{xx}v_\eta.
\end{aligned}
\]


Làm tương tự cho $u_{xy},u_{yy}$. Có thể dùng toán tử $\partial_x,\partial_y$ để rút ngắn, nhưng phải nêu các đạo hàm của hệ số nếu đổi biến phi tuyến.

\textbf{Bước 4 — Gom hệ số và chia hợp lệ.} Hệ số mới của $v_{\xi\xi}$ là $A\xi_x^2+2B\xi_x\xi_y+C\xi_y^2$; hệ số của $v_{\xi\eta}$ là $2[A\xi_x\eta_x+B(\xi_x\eta_y+\xi_y\eta_x)+C\xi_y\eta_y]$. Hệ số mới của $v_\xi$ là $A\xi_{xx}+2B\xi_{xy}+C\xi_{yy}+D\xi_x+E\xi_y$. Chỉ chia cho hệ số khác $0$ trên miền đang xét.

\textbf{Bước 5 — Với ba biến.} Lập ma trận đối xứng $M$ của phần cấp hai; hệ số hỗn hợp được chia đôi. Dùng hoàn thành bình phương để tìm số hướng dương, âm và bằng $0$. Nếu $\xi=T x$, ma trận mới là $TMT^T$. Cho rõ ma trận $T$ và kiểm tra nó khả nghịch; chỉ ghi dấu của dạng toàn phương chưa đủ để thực hiện phép đổi biến.

\textbf{Chi tiết gom hệ số sau đổi biến.} Với $Lu=Au_{xx}+2Bu_{xy}+Cu_{yy}+Du_x+Eu_y+Ku$, tính riêng


\[
\begin{aligned}
a&=A\xi_x^2+2B\xi_x\xi_y+C\xi_y^2,\\
2b&=2[A\xi_x\eta_x+B(\xi_x\eta_y+\xi_y\eta_x)+C\xi_y\eta_y],\\
c&=A\eta_x^2+2B\eta_x\eta_y+C\eta_y^2,\\
d&=A\xi_{xx}+2B\xi_{xy}+C\xi_{yy}+D\xi_x+E\xi_y,\\
e&=A\eta_{xx}+2B\eta_{xy}+C\eta_{yy}+D\eta_x+E\eta_y.
\end{aligned}
\]


Khi đó $Lu=av_{\xi\xi}+2bv_{\xi\eta}+cv_{\eta\eta}+dv_\xi+ev_\eta+Kv$. Trong lời giải từng ý, thế các số hoặc hàm hệ số vào từng dòng trước khi rút gọn. Với ba biến, tính $TM$ trước, rồi lấy tích vô hướng mỗi hàng của $TM$ với mỗi hàng của $T$ để có $TMT^T$.
\subsection*{Lời giải áp dụng cho bài}
\textbf{Bước 1. Phân loại.} Phương trình là $2u_{xx}+3u_{xy}+u_{yy}+u_x-u_y=0$. Lấy $A=2$, $B=3/2$, $C=1$, ta được $\delta=9/4-2=1/4>0$: hyperbolic trên $\mathbb R^2$.

\textbf{Bước 2. Hai họ đặc trưng.} Phương trình đặc trưng
\[2(dy)^2-3dxdy+(dx)^2=(dy-dx)(2dy-dx)=0\]
cho $y-x=\mathrm{const}$ và $2y-x=\mathrm{const}$. Chọn $\xi=x-y$, $\eta=x-2y$. Jacobian $J=-1\ne0$; phép đổi ngược là $x=2\xi-\eta$, $y=\xi-\eta$.

\textbf{Bước 3. Thay đạo hàm.} Vì $\partial_x=\partial_\xi+\partial_\eta$ và $\partial_y=-\partial_\xi-2\partial_\eta$, ta có
\begin{align*}
u_{xx}&=v_{\xi\xi}+2v_{\xi\eta}+v_{\eta\eta},\\
u_{xy}&=-v_{\xi\xi}-3v_{\xi\eta}-2v_{\eta\eta},\\
u_{yy}&=v_{\xi\xi}+4v_{\xi\eta}+4v_{\eta\eta},\\
u_x-u_y&=2v_\xi+3v_\eta.
\end{align*}
Các hệ số của $v_{\xi\xi}$, $v_{\eta\eta}$ lần lượt là $2-3+1=0$, $2-6+4=0$; hệ số hỗn hợp là $4-9+4=-1$.

\textbf{Kết quả.} Dạng chính tắc là
\[\boxed{-v_{\xi\eta}+2v_\xi+3v_\eta=0.}\]
Không bỏ các hạng cấp một khi phân loại phần chính. Có thể nhân cả phương trình với $-1$ để viết $v_{\xi\eta}-2v_\xi-3v_\eta=0$.\end{document}